
\textbf{Re-state the overall trend.} The polarization morphology of IRS 63 was presented in Chapter \ref{sec: c4_polarization_maps}, and showed that it has a clear wavelength dependence. The morphology transitioned from mostly azimuthal alignment at \bfour to mostly aligned with the disk minor axis at \bsevenNS. Between the upper and lower bounds of observing wavelength, the morphology was intermediate, or a combination of the two. These morphologies, especially the intermediate cases, indicate that the morphology, caused by polarization mechanisms,s cannot be described by a single mechanism, and there must be multiple mechanisms contributing to the observed emission. 


\textbf{Morphology to Mechanism with Wavelength.} The azimuthal pattern observed mostly at longer wavelengths is consistent with the pattern expected from aligned dust grains (e.g., via radiative alignment or aerodynamic alignment, see Chapter \ref{sec: intro_polarization_mechanisms}), as shown in Figure \ref{fig: kataoka_polarization_morphologies}(b). The minor-axis aligned morphology is consistent with dust self-scattering, as shown in Figure \ref{fig: kataoka_polarization_morphologies}(c). This indicates that dust self-scattering contributes more to the polarization emission at shorter wavelengths. 


\textbf{Central vs Outer Disk.} Even at \bfourNS, where the azimuthal pattern dominates the morphology, the vectors are still aligned with the disk minor axis at the very centre. Alternatively, at the lowest wavelength of \bsevenNS, where we see the most uniform alignment, there is still an azimuthal pattern at the outer edge of the disk. This trend is also seen at the intermediate wavelengths. This indicates that not only is the dominant mechanism changing with wavelength, but at a single observing wavelength, the dominant polarization mechanism varies across the disk as well. \note{at all wavelengths, centre and outer look different}


\textbf{Spatial changes w.r.t beam size.} \add{Area in AU based on beam, band 4 small but smallest beam}




\textbf{Optical depth.} The measurement of optical depth refers to how easily photons can pass through a medium \citep{Kuruwita}. When something is optically thin, it has a lower optical depth, and light passes through easily \citep{Kuruwita, Lin2022}. At optically thick, high optical depths, grains encounter other grains more often and are absorbed, re-emitted and scattered \citep{Kuruwita, Lin2022}. The optical depth of the continuum emission increases as wavelength decreases \citep{Williams}, and dust self-scattering tends to dominate at higher optical depths \citep{Lin2022}. This is consistent with the observations, as we see more signatures of dust self-scattering at higher optical depth (lower wavelength). As the disk's density in dust decreases radially outward, so does the optical depth \citep{Williams}. This also makes sense as we see less dust self-scattering in the optically thin outer regions of the disk.

Overall, this transition in polarization morphology both with wavelength and position across the disk suggests that the dominant polarization mechanism changes not only with wavelength but also spatially across the disk. \add{connect to models for next section?}
%Thermal polarization dominates at low optical depths \citep{Lin2022}. 
=

% Young disks are optically thick \citep{Kuruwita}. 

